\documentclass[12pt,a4paper]{exam}
\usepackage[utf8]{inputenc}
\usepackage{amsmath,amssymb,amsfonts}
\usepackage{geometry}
\usepackage{enumitem}
\usepackage{graphicx}
\usepackage{hyperref}
\usepackage{xcolor}
\geometry{a4paper, top=2cm, bottom=2cm, left=2.5cm, right=2.5cm}
\hypersetup{colorlinks=true, linkcolor=blue, urlcolor=blue}
\renewcommand{\thequestion}{\textbf{\arabic{question}}}
\renewcommand{\questionlabel}{\thequestion.}
\pointname{ marks}
\pointsdroppedatright
\bracketedpoints
\setlength{\parindent}{0pt}
\setlength{\emergencystretch}{3em}
\allowdisplaybreaks
\newsavebox{\schmathbox}
\newcommand{\fitmath}[1]{%
  \sbox{\schmathbox}{$\displaystyle #1$}%
  \ifdim\wd\schmathbox>\linewidth
    \resizebox{\linewidth}{!}{\usebox{\schmathbox}}%
  \else
    \usebox{\schmathbox}%
  \fi}

\begin{document}

\thispagestyle{empty}
\begin{center}
  \textbf{\LARGE Diag} \\[0.3cm]
  \textbf{Time Allowed: 3 Hours} \hfill \textbf{Maximum Marks: 80} \\[0.3cm]
  \rule{\textwidth}{0.5pt}
\end{center}

\vspace{0.3cm}

\noindent\textbf{General Instructions:}
\begin{enumerate}
  \item {\normalfont\normalcolor This question paper contains 40 questions. All questions are compulsory.}
  \item {\normalfont\normalcolor Questions are grouped into sections A through E.}
  \item {\normalfont\normalcolor Use of calculators is NOT permitted.}
\end{enumerate}
\bigskip
\noindent\textbf{\large Section A: Multiple Choice \& Assertion-Reason (1 mark each)}

\begin{questions}
\question[1] If a tangent $PA$ and a tangent $PB$ from a point $P$ to a circle with centre $O$ are inclined to each other at an angle of $80^\circ$, then $\angle POA$ is equal to:

\begin{center}
\includegraphics[width=0.48\textwidth]{/home/user/Desktop/backup/schEngine/pdfs/diag/diagrams/diagram_047.pdf}
\end{center}

\begin{enumerate}[label=(\Alph*)]
  \item $(A) 50^\circ$
  \item $(B) 60^\circ$
  \item $(C) 70^\circ$
  \item $(D) 80^\circ$
\end{enumerate}
\vspace{0.15cm}

\question[1] A circle touches all four sides of a quadrilateral $\text{ABCD}$. If $AB = 6 \text{ cm}$, $BC = 7 \text{ cm}$, and $CD = 4 \text{ cm}$, then the length of $AD$ is:
\begin{enumerate}[label=(\Alph*)]
  \item $(A) 2 \text{ cm}$
  \item $(B) 2.5 \text{ cm}$
  \item $(C) 3 \text{ cm}$
  \item $(D) 3.5 \text{ cm}$
\end{enumerate}
\vspace{0.15cm}

\question[1] A tangent $PQ$ at a point $P$ of a circle of radius $5 \text{ cm}$ meets a line through the centre $O$ at a point $Q$ so that $OQ = 13 \text{ cm}$. The length of $PQ$ is:

\begin{center}
\includegraphics[width=0.48\textwidth]{/home/user/Desktop/backup/schEngine/pdfs/diag/diagrams/diagram_048.pdf}
\end{center}

\begin{enumerate}[label=(\Alph*)]
  \item (A) $10 \text{ cm}$
  \item (B) $11 \text{ cm}$
  \item (C) $12 \text{ cm}$
  \item (D) $8 \text{ cm}$
\end{enumerate}
\vspace{0.15cm}

\question[1] If tangents $PA$ and $PB$ from a point $P$ to a circle with centre $O$ are inclined to each other at an angle of $80^\circ$, then $\angle POA$ is equal to:
\begin{enumerate}[label=(\Alph*)]
  \item (A) $50^\circ$
  \item (B) $60^\circ$
  \item (C) $70^\circ$
  \item (D) $80^\circ$
\end{enumerate}
\vspace{0.15cm}

\question[1] The length of a tangent from a point $A$ at a distance of $5 \text{ cm}$ from the centre of the circle is $4 \text{ cm}$. The radius of the circle is:
\begin{enumerate}[label=(\Alph*)]
  \item (A) $2 \text{ cm}$
  \item (B) $3 \text{ cm}$
  \item (C) $4 \text{ cm}$
  \item (D) $5 \text{ cm}$
\end{enumerate}
\vspace{0.15cm}

\question[1] Two concentric circles are of radii $5 \text{ cm}$ and $3 \text{ cm}$. The length of the chord of the larger circle which touches the smaller circle is:
\begin{enumerate}[label=(\Alph*)]
  \item (A) $6 \text{ cm}$
  \item (B) $7 \text{ cm}$
  \item (C) $8 \text{ cm}$
  \item (D) $10 \text{ cm}$
\end{enumerate}
\vspace{0.15cm}

\question[1] A quadrilateral $\text{ABCD}$ is drawn to circumscribe a circle. If $AB = 6 \text{ cm}$, $BC = 7 \text{ cm}$, and $CD = 4 \text{ cm}$, then $AD$ is equal to:
\begin{enumerate}[label=(\Alph*)]
  \item (A) $2 \text{ cm}$
  \item (B) $4 \text{ cm}$
  \item (C) $5 \text{ cm}$
  \item (D) $3 \text{ cm}$
\end{enumerate}
\vspace{0.15cm}

\question[1] The distance of the point $P(3, 4)$ from the origin $(0, 0)$ is:

\begin{center}
\includegraphics[width=0.48\textwidth]{/home/user/Desktop/backup/schEngine/pdfs/diag/diagrams/diagram_049.pdf}
\end{center}

\begin{enumerate}[label=(\Alph*)]
  \item (A) $3$
  \item (B) $4$
  \item (C) $5$
  \item (D) $7$
\end{enumerate}
\vspace{0.15cm}

\question[1] The distance of the point $P(3, 4)$ from the origin is
\begin{enumerate}[label=(\Alph*)]
  \item (A) $3$
  \item (B) $5$
  \item (C) $7$
  \item (D) $25$
\end{enumerate}
\vspace{0.15cm}

\question[1] The coordinates of the midpoint of the line segment joining the points $A(-2, 8)$ and $B(-6, -4)$ are
\begin{enumerate}[label=(\Alph*)]
  \item (A) $(-4, -6)$
  \item (B) $(2, 6)$
  \item (C) $(-4, 2)$
  \item (D) $(-8, 4)$
\end{enumerate}
\vspace{0.15cm}

\question[1] The point on the $x$-axis which is equidistant from points $A(2, -5)$ and $B(-2, 9)$ is
\begin{enumerate}[label=(\Alph*)]
  \item (A) $(-7, 0)$
  \item (B) $(0, -7)$
  \item (C) $(7, 0)$
  \item (D) $(0, 7)$
\end{enumerate}
\vspace{0.15cm}

\question[1] If the point $P(k, 0)$ divides the line segment joining $A(2, -2)$ and $B(-7, 4)$ in the ratio $1:2$, the value of $k$ is
\begin{enumerate}[label=(\Alph*)]
  \item (A) $2$
  \item (B) $-2$
  \item (C) $1$
  \item (D) $-1$
\end{enumerate}
\vspace{0.15cm}

\question[1] The distance between the points $P(a, 0)$ and $Q(0, b)$ is
\begin{enumerate}[label=(\Alph*)]
  \item (A) $a + b$
  \item (B) $\sqrt{a^2 + b^2}$
  \item (C) $a^2 + b^2$
  \item (D) $\sqrt{a+b}$
\end{enumerate}
\vspace{0.15cm}

\question[1] Assertion (A): The length of a tangent drawn from an external point to a circle is always equal to the radius of the circle. Reason (R): A tangent to a circle is perpendicular to the radius through the point of contact.

\begin{center}
\includegraphics[width=0.48\textwidth]{/home/user/Desktop/backup/schEngine/pdfs/diag/diagrams/diagram_050.pdf}
\end{center}

\begin{enumerate}[label=(\Alph*)]
  \item (A) Both Assertion (A) and Reason (R) are true and Reason (R) is the correct explanation of Assertion (A)
  \item (B) Both Assertion (A) and Reason (R) are true but Reason (R) is NOT the correct explanation of Assertion (A)
  \item (C) Assertion (A) is true but Reason (R) is false
  \item (D) Assertion (A) is false but Reason (R) is true
\end{enumerate}
\vspace{0.15cm}

\question[1] Assertion (A): The number of tangents that can be drawn from a point inside a circle is zero. Reason (R): A tangent to a circle intersects the circle at only one point.
\begin{enumerate}[label=(\Alph*)]
  \item (A) Both Assertion (A) and Reason (R) are true and Reason (R) is the correct explanation of Assertion (A)
  \item (B) Both Assertion (A) and Reason (R) are true but Reason (R) is NOT the correct explanation of Assertion (A)
  \item (C) Assertion (A) is true but Reason (R) is false
  \item (D) Assertion (A) is false but Reason (R) is true
\end{enumerate}
\vspace{0.15cm}

\question[1] Assertion (A): If the angle between two tangents drawn from a point $P$ to a circle of radius $r$ and center $O$ is $60^\circ$, then $OP = 2r$. Reason (R): The line segment joining the center to the external point bisects the angle between the two tangents.
\begin{enumerate}[label=(\Alph*)]
  \item (A) Both Assertion (A) and Reason (R) are true and Reason (R) is the correct explanation of Assertion (A)
  \item (B) Both Assertion (A) and Reason (R) are true but Reason (R) is NOT the correct explanation of Assertion (A)
  \item (C) Assertion (A) is true but Reason (R) is false
  \item (D) Assertion (A) is false but Reason (R) is true
\end{enumerate}
\vspace{0.15cm}

\question[1] Assertion (A): Two concentric circles have the same center. Reason (R): A chord of the larger circle which touches the smaller circle is bisected at the point of contact.
\begin{enumerate}[label=(\Alph*)]
  \item (A) Both Assertion (A) and Reason (R) are true and Reason (R) is the correct explanation of Assertion (A)
  \item (B) Both Assertion (A) and Reason (R) are true but Reason (R) is NOT the correct explanation of Assertion (A)
  \item (C) Assertion (A) is true but Reason (R) is false
  \item (D) Assertion (A) is false but Reason (R) is true
\end{enumerate}
\vspace{0.15cm}

\question[1] Assertion (A): The length of the tangent from a point $A$ at a distance of $5 \text{ cm}$ from the centre of the circle of radius $3 \text{ cm}$ is $4 \text{ cm}$. Reason (R): The tangent at any point of a circle is perpendicular to the radius through the point of contact.

\begin{center}
\includegraphics[width=0.48\textwidth]{/home/user/Desktop/backup/schEngine/pdfs/diag/diagrams/diagram_051.pdf}
\end{center}

\begin{enumerate}[label=(\Alph*)]
  \item (A) Both Assertion (A) and Reason (R) are true and Reason (R) is the correct explanation of Assertion (A)
  \item (B) Both Assertion (A) and Reason (R) are true but Reason (R) is NOT the correct explanation of Assertion (A)
  \item (C) Assertion (A) is true but Reason (R) is false
  \item (D) Assertion (A) is false but Reason (R) is true
\end{enumerate}
\vspace{0.15cm}

\question[1] Assertion (A): The distance of the point $P(3, 4)$ from the origin is $5$ units. Reason (R): The distance of a point $(x, y)$ from the origin $(0, 0)$ is given by $\sqrt{x^2 + y^2}$.
\begin{enumerate}[label=(\Alph*)]
  \item (A) Both Assertion (A) and Reason (R) are true and Reason (R) is the correct explanation of Assertion (A)
  \item (B) Both Assertion (A) and Reason (R) are true but Reason (R) is NOT the correct explanation of Assertion (A)
  \item (C) Assertion (A) is true but Reason (R) is false
  \item (D) Assertion (A) is false but Reason (R) is true
\end{enumerate}
\vspace{0.15cm}

\question[1] Assertion (A): The midpoint of the line segment joining $A(2, 4)$ and $B(4, 2)$ is $(3, 3)$. Reason (R): The midpoint of a line segment with endpoints $(x_1, y_1)$ and $(x_2, y_2)$ is $(\frac{x_1 + x_2}{2}, \frac{y_1 + y_2}{2})$.
\begin{enumerate}[label=(\Alph*)]
  \item (A) Both Assertion (A) and Reason (R) are true and Reason (R) is the correct explanation of Assertion (A)
  \item (B) Both Assertion (A) and Reason (R) are true but Reason (R) is NOT the correct explanation of Assertion (A)
  \item (C) Assertion (A) is true but Reason (R) is false
  \item (D) Assertion (A) is false but Reason (R) is true
\end{enumerate}
\vspace{0.15cm}

\question[1] Assertion (A): The point $(0, 5)$ lies on the $y$-axis. Reason (R): For any point on the $y$-axis, the $x$-coordinate is zero.
\begin{enumerate}[label=(\Alph*)]
  \item (A) Both Assertion (A) and Reason (R) are true and Reason (R) is the correct explanation of Assertion (A)
  \item (B) Both Assertion (A) and Reason (R) are true but Reason (R) is NOT the correct explanation of Assertion (A)
  \item (C) Assertion (A) is true but Reason (R) is false
  \item (D) Assertion (A) is false but Reason (R) is true
\end{enumerate}
\vspace{0.15cm}

\question[1] Assertion (A): The distance between the points $A(0, 0)$ and $B(0, 7)$ is $7$ units. Reason (R): The distance between two points $(x_1, y_1)$ and $(x_2, y_2)$ is $\sqrt{(x_2-x_1)^2 + (y_2-y_1)^2}$.
\begin{enumerate}[label=(\Alph*)]
  \item (A) Both Assertion (A) and Reason (R) are true and Reason (R) is the correct explanation of Assertion (A)
  \item (B) Both Assertion (A) and Reason (R) are true but Reason (R) is NOT the correct explanation of Assertion (A)
  \item (C) Assertion (A) is true but Reason (R) is false
  \item (D) Assertion (A) is false but Reason (R) is true
\end{enumerate}
\vspace{0.15cm}

\question[1] Assertion (A): The point $(3, 0)$ lies on the $x$-axis. Reason (R): For any point on the $x$-axis, the $y$-coordinate is zero.
\begin{enumerate}[label=(\Alph*)]
  \item (A) Both Assertion (A) and Reason (R) are true and Reason (R) is the correct explanation of Assertion (A)
  \item (B) Both Assertion (A) and Reason (R) are true but Reason (R) is NOT the correct explanation of Assertion (A)
  \item (C) Assertion (A) is true but Reason (R) is false
  \item (D) Assertion (A) is false but Reason (R) is true
\end{enumerate}
\vspace{0.15cm}

\end{questions}
\bigskip
\noindent\textbf{\large Section B: Very Short Answer (2 marks each)}

\begin{questions}
\question[2] In a circle with centre $O$, a tangent $PT$ is drawn from a point $P$ outside the circle. If $OP = 13$ $cm$ and the radius of the circle is $5$ $cm$, find the length of the tangent $PT$.

\begin{center}
\includegraphics[width=0.48\textwidth]{/home/user/Desktop/backup/schEngine/pdfs/diag/diagrams/diagram_052.pdf}
\end{center}

\vspace{0.15cm}

\question[2] Two concentric circles are of radii $5$ $cm$ and $3$ $cm$. Find the length of the chord of the larger circle which touches the smaller circle.

\begin{center}
\includegraphics[width=0.48\textwidth]{/home/user/Desktop/backup/schEngine/pdfs/diag/diagrams/diagram_053.pdf}
\end{center}

\vspace{0.15cm}

\question[2] $PA$ and $PB$ are two tangents from a point $P$ to a circle with centre $O$ such that $\angle APB = 80^\circ$. Find $\angle AOB$.
\vspace{0.15cm}

\question[2] If a tangent $PA$ and $PB$ from a point $P$ to a circle with centre $O$ are inclined to each other at an angle of $70^\circ$, then find $\angle POA$.
\vspace{0.15cm}

\question[2] A point $P$ is $26$ $cm$ away from the centre $O$ of a circle and the length of the tangent $PT$ drawn from $P$ to the circle is $10$ $cm$. Find the radius of the circle.
\vspace{0.15cm}

\end{questions}
\bigskip
\noindent\textbf{\large Section C: Short Answer (3 marks each)}

\begin{questions}
\question[3] A tangent $PQ$ at a point $P$ of a circle of radius $5$ $cm$ meets a line through the centre $O$ at a point $Q$ so that $OQ = 13$ $cm$. Find the length of $PQ$.

\begin{center}
\includegraphics[width=0.48\textwidth]{/home/user/Desktop/backup/schEngine/pdfs/diag/diagrams/diagram_054.pdf}
\end{center}

\vspace{0.15cm}

\question[3] Two concentric circles are of radii $5$ $cm$ and $3$ $cm$. Find the length of the chord of the larger circle which touches the smaller circle.

\begin{center}
\includegraphics[width=0.48\textwidth]{/home/user/Desktop/backup/schEngine/pdfs/diag/diagrams/diagram_055.pdf}
\end{center}

\vspace{0.15cm}

\question[3] If tangents $PA$ and $PB$ from a point $P$ to a circle with centre $O$ are inclined to each other at angle of $80^\circ$, then find $\angle POA$.

\begin{center}
\includegraphics[width=0.48\textwidth]{/home/user/Desktop/backup/schEngine/pdfs/diag/diagrams/diagram_056.pdf}
\end{center}

\vspace{0.15cm}

\question[3] Prove that the tangents drawn at the ends of a diameter of a circle are parallel.
\vspace{0.15cm}

\question[3] A quadrilateral $\text{ABCD}$ is drawn to circumscribe a circle. Prove that $AB + CD = AD + BC$.
\vspace{0.15cm}

\question[3] Find the ratio in which the $y$-axis divides the line segment joining the points $A(5, -6)$ and $B(-1, -4)$. Also, find the point of intersection.

\begin{center}
\includegraphics[width=0.48\textwidth]{/home/user/Desktop/backup/schEngine/pdfs/diag/diagrams/diagram_057.pdf}
\end{center}

\vspace{0.15cm}

\end{questions}
\bigskip
\noindent\textbf{\large Section D: Long Answer (4-5 marks each)}

\begin{questions}
\question[5] Prove that the lengths of tangents drawn from an external point to a circle are equal. Using this result, find the length of $PQ$ if $PQ$ is a tangent to a circle with centre $O$ and radius $5 \text{ cm}$, where $OQ = 13 \text{ cm}$ and $P$ is the external point.

\begin{center}
\includegraphics[width=0.48\textwidth]{/home/user/Desktop/backup/schEngine/pdfs/diag/diagrams/diagram_058.pdf}
\end{center}

\vspace{0.15cm}

\question[5] A quadrilateral $\text{ABCD}$ is drawn to circumscribe a circle. Prove that $AB + CD = AD + BC$.
\vspace{0.15cm}

\question[5] Two concentric circles are of radii $5 \text{ cm}$ and $3 \text{ cm}$. Find the length of the chord of the larger circle which touches the smaller circle.
\vspace{0.15cm}

\question[5] If a triangle $ABC$ is drawn to circumscribe a circle of radius $2 \text{ cm}$ such that the segments $BD$ and $DC$ into which $BC$ is divided by the point of contact $D$ are of lengths $4 \text{ cm}$ and $3 \text{ cm}$ respectively, find the sides $AB$ and $AC$.
\vspace{0.15cm}

\question[4] Find the ratio in which the $y$-axis divides the line segment joining the points $A(5, -6)$ and $B(-1, -4)$. Also, find the point of intersection.

\begin{center}
\includegraphics[width=0.48\textwidth]{/home/user/Desktop/backup/schEngine/pdfs/diag/diagrams/diagram_061.pdf}
\end{center}

\vspace{0.15cm}

\question[4] The vertices of $\text{riangle ABC}$ are $A(4, 6)$, $B(1, 5)$, and $C(7, 2)$. A line is drawn to intersect sides $AB$ and $AC$ at points $D$ and $E$ respectively, such that $\frac{AD}{AB} = \frac{AE}{AC} = \frac{1}{4}$. Find the coordinates of points $D$ and $E$ and show that the area of $\text{riangle ADE}$ is $\frac{1}{16}$ of the area of $\text{riangle ABC}$.

\begin{center}
\includegraphics[width=0.48\textwidth]{/home/user/Desktop/backup/schEngine/pdfs/diag/diagrams/diagram_062.pdf}
\end{center}

\vspace{0.15cm}

\end{questions}
\bigskip

\end{document}
